% !TEX TS-program = lualatex+se
% !TEX encoding = UTF-8 Unicode

\documentclass[14pt, a4paper]{extbook}

% --- USTAWIENIA STRONY ---
\usepackage{geometry}
\geometry{
	left=15mm,
	right=15mm,
	top=15mm,
	bottom=15mm,
	marginparwidth=0mm,
	marginparsep=0mm
	}
%\usepackage[letterpaper,showframe]{geometry}	
% usuwanie nagłówków/stopek z\,dokumentu
\usepackage[indent=0pt]{parskip} % ustawianie paragrafów
%\setstretch{1.5} % półtorej linii odległości między wierszami
\usepackage{titlesec} % ustawianie nagłówków
%\usepackage{fullpage} % redukcja marginesów
\usepackage{setspace} % odstępy między liniami
\usepackage{enumitem}% listy numerowane
\usepackage{soul}
% -----------------------------
\renewcommand*{\thepage}{\footnotesize{\arabic{page}}}% mniejsza czcionka numeru strony, dodanie czegokolwiek przed numerem: {A--\arabic{page}}
\newcommand{\margintext}[1]{\textit{\small #1}}
% przesunięcia poziome
\newcommand{\smallwidth}{\hskip 1em}
\newcommand{\medwidth}{\hskip 3pt}
\newcommand{\largewidth}{\hskip 12pt}
% przesunięcia pionowe
\newcommand{\linespace}{\vskip 7pt}
\newcommand{\paragraphspace}{\vskip 8pt}
\newcommand{\segmentspace}{\vskip 18pt}
\newcommand{\sectionspace}{\vskip 28pt}
%------------------------------------------
\newcommand{\wstawplik}[1]{\input{/Users/Laser/Documents/Latex/include/#1}}%skraca kod
% --- GREGORIO ---	
\usepackage{gregoriotex}
\gresetlinecolor{gregoriocolor}% Make staff lines red; remove for black:
%\gresetgregoriofont{granapadano} %{gregorio}
\gresetheadercapture{annotation}{greannotation}{string}% odczytuje adnotacje nad inicjałem jako ciąg, co pozwala na jego renderowanie
\gresetheadercapture{commentary}{grecommentary}{string}% odczytuje komentarz jako ciąg, co pozwala na jego renderowanie
\grechangedim{commentaryraise}{0.3cm}{scalable} % podnosi komentarz wystarczająco ponad górną linię
\grechangestyle{annotation}{\color{black}\small}% format adnotacji
\grechangestyle{commentary}{\color{gregoriocolor}\small}% format komentarza
\grechangestyle{abovelinestext}{\color{gregoriocolor}\small}% format napisu nad liniami
\grechangestyle{initial}{\fontsize{36}{36}\selectfont} % zmniejsza wielkość inicjału do bardziej proporcjonalnego rozmiaru (=36pt)
\grechangedim{intersyllablespacenotes}{0.2cm}{scalable}% Minimalna odległość między nutami różnych sylab.
% Set the space around the initial:
% See http://gregorio-project.github.io/gregoriotex/details.html for more details and options
\grechangedim{beforeinitialshift}{0.5mm}{scalable}
\grechangedim{afterinitialshift}{1.0mm}{scalable}

% umieszczenie zawartości składu ich unikalnego argumentu po lewej i prawej stronie każdej linii
\let\grelocalleftbox\localleftbox
\let\grelocalrightbox\localrightbox

%\newcommand{\separator}{\greseparator{3}{40}}
%\newfontface\vectis{VectisMonumentalW01-Regular}
%\def\greinitialformat#1{\fontsize{20}{20}\vectis #1}
% ------------------

\usepackage[hyperindex=true]{hyperref}

%---- usun konflikt między paracol i luacolor ----
%\wstawplik{conparlua}
%-----------------------------------------------

% --- SPIS TREŚCI -------
%\usepackage{imakeidx}
%\makeindex[title=Alfabetyczny spis program=texindy,options=-M styl.xdy]
%\usepackage{filecontents}
%\begin{filecontents*}{styl.xdy}
%;;; xindy style file
%(markup-locclass-list :open "\dotfill" :sep "")
%(define-attributes (( "textbf" "default" )) )
%(markup-locref   :attr  "textbf"     :open "\textbf{\hyperpage{" :close "}}")
%(markup-locref   :attr  "textit"     :open "\textit{\hyperpage{" :close "}}")
%(markup-locref   :attr  "textttt"     :open "\textttt{\hyperpage{" :close "}}")
%(markup-locref   :attr  "texttsc"     :open "\texttsc{\hyperpage{" :close "}}")
%(markup-locref   :attr  "default"     :open "\hyperpage{" :close "}")
%\end{filecontents*}
% -------------------------

% ---  Konfiguracja kolumn ---
\usepackage{paracol}
\setcolumnwidth{.49\textwidth, .49\textwidth}
\setlength{\columnseprule}{0.4pt}
\setlength{\columnsep}{.02\textwidth}
\colseprulecolor{red}% wymagany pakiet xcolor
\footnotelayout{m}
% ----------------------------

%\makeatletter 
%\patchcmd\pcol@hfil
% {\copy\pcol@tempboxa}
% {\luacolorresetboxcolor\pcol@tempboxa\copy\pcol@tempboxa}{}{\fail}
%\makeatother
% ----------------------------

\footnotelayout{m}

%% --- UKŁAD DWUKOLUMNOWY Z TŁUMACZENIEM ---
%\usepackage{tcolorbox}
%\newtcolorbox{nutybox}{size=minimal, colback=white, colframe=red, arc=0mm, width=\linewidth, halign=left, valign=center}
%\newtcolorbox{tlumbox}{size=minimal, colback=white, colframe=red, arc=0mm, width=\linewidth, halign=left, valign=top, fontupper=\small\itshape} 
%\newcommand{\nuty}[1]{\begin{nutybox}#1 \end{nutybox}}
%\newcommand{\tlum}[1]{\begin{tlumbox}\tlumacz{#1}\end{tlumbox}}
%% ------------------------------------------------------

\usepackage{Tabbing}
\usepackage{longtable}
\usepackage[savepos]{zref}

% --- CZCIONKI DOKUMENTU ---
\usepackage[latin,polish]{babel}
\usepackage{fontspec}
%\usepackage{times}
%\usepackage[math]{anttor}
%\usepackage{tgtermes}
%\usepackage{mathptmx}
\usepackage{ebgaramond}
\usepackage[osf,p]{libertine}
\defaultfontfeatures{Ligatures=TeX}

%\setmainfont{TeX Gyre Termes}
%\setmainfont{TeX Gyre Heros}
%\setmainfont{Calluna Regular}
\setmainfont{EB Garamond}
%\setmainfont{Bodoni Moda}
%\newfontfamily{\sekfont}{Times New Roman}
\newfontfamily{\sekfont}{Libertinus Serif}
\newcommand\sekcjafont{\fontsize{16}{16}\selectfont \sekfont}

%----- Strona tytułowa -------------
\newfontfamily{\tytfont}{JSL Ancient}
\newcommand\tytulfont[1]{\fontsize{36}{30}\selectfont \tytfont#1\par}
\newcommand\podtytulfont[1]{\fontsize{18}{10}\selectfont \tytfont#1\par}

%\setmainfont[<font features>]{<name of font>} % sets the roman font
%\setsansfont[<font features>]{<name of font>} % sets the sans font
%\setmonofont[<font features>]{<name of font>} % sets the monospace font

% --- KONIEC CZCIONKI DOKUMENTU ---

% --- GRAFIKA I KOLORY ---
\usepackage{xcolor} % kolory
\usepackage{graphicx} % umożliwia import obrazów
\graphicspath{{img/}} % katalog obrazów
\usepackage{pgfornament}% ornamenty graficzne
% --- KONIEC GRAFIKA I KOLORY -----------

% --- KOMENDY WŁASNE ---
\newcommand{\red}[1]{\textcolor{red}{#1}}
\newcommand{\black}[1]{\textcolor{black}{#1}}
% --------------------------------
%\newcommand{\firstchar}[1]{\directlua{tex.sprint(unicode.utf8.sub('#1',1,1))}}%makro pobiera pierwszy znak ze stringu
\newcommand{\tlumacz}[1]{\small\textit{#1}}
\newcommand{\inicjal}[2]{\textcolor{red}{\Large\textbf{#1}}{#2}}
\newcommand{\tlumaczenie}[3]{\selectlanguage{polish}\textcolor{red}{\Large\textbf{#1}}{#2}{#3}}%
%--------------------------------
% ----------- RUBRYKA -------------------
\newcommand{\rubryka}[1]{
		\parbox{\textwidth}{\hspace{5mm}\small\textit{\red{#1}}}
	\linespace}
%-----------STYLE NAGŁÓWKÓW ----------------
%
% ----------- ROZDZIAŁ --------------------
\titlespacing*{\chapter}{0pt}{-40pt}{0pt}	
\titleformat{\chapter}[hang]{\filcenter\scshape\bfseries\Large}{}{0pt}{} % format nagłówków
\newcommand{\rozdzial}[2]{
		\chapter{#1}
		\vspace{-1em}
		\begin{spacing}{0.8}
			\begin{center}
				\parbox{0.9\textwidth}{\centering\small\textit{\red{#2}}}
			\end{center}
		\end{spacing}
		\linespace}

% ----------- SEKCJA --------------------
\titlespacing*{\section}{0pt}{0pt}{0pt}	
\titleformat{\section}[hang]{\filcenter\scshape\normalfont}{}{0pt}{} % format nagłówków
\newcommand{\sekcja}[2]{
	\begin{spacing}{0.6}
	\begin{center}
		\segmentspace
		\section*{#1}
		\begin{minipage}[c]{0.88\linewidth}
			\rubryka{#2}
			%\centering\red{\small\textit{#2}}
		\end{minipage}
	\end{center}	
	\end{spacing}
	\linespace}
	
% ----------- SUBSEKCJA --------------------	
\titlespacing*{\subsection}{0pt}{0pt}{0pt}	
\titleformat{\subsection}[hang]{\normalfont\scshape}{}{0pt}{} % format podnagłówków
\newcommand{\subsekcja}[1]{
	\segmentspace
	\subsection*{\red{#1}}		
	\segmentspace}
%------------ Tytuł -----------------------
\newcommand{\tytul}[1]{
	\red{\large\scshape{#1}\linespace}}
%------------ Podytuł -----------------------
\newcommand{\podtytul}[1]{
	\red{\normalfont\scshape{#1}\linespace}}
%------------DEKOR ---------------------
\newcommand{\dekor}[2]{
	\begin{center}
		\vspace*{\fill}
		\includegraphics[width=#1\textwidth]{#2}
		\vspace*{\fill}
	\end{center}}
%------- Zarządzanie wersami w psalmach z tlumaczeniem -------------------------
\newcounter{licznik}
\newcommand{\licznikZero}{\setcounter{licznik}{0}}% zeruj licznik
\newcommand{\licznikUp}{\refstepcounter{licznik}}% zwiększ licznik 
\licznikZero
\newcommand{\BeginParallel}{\begin{spacing}{0.8}\begin{paracol}{2}\licznikZero}
\newcommand{\lat}[1]{\selectlanguage{latin}\ensurevspace{2\baselineskip}\begin{minipage}[t]{\linewidth}\ifnum \value{licznik}>0 \licznikUp \thelicznik .~#1 \else \licznikUp#1 \fi \end{minipage}\switchcolumn}
\newcommand{\latN}[1]{\selectlanguage{latin}\ensurevspace{2\baselineskip}\begin{minipage}[t]{\linewidth}#1\end{minipage}\switchcolumn}% kolumna bez numeracji wersów
\newcommand{\pol}[1]{\selectlanguage{polish}\begin{minipage}[t]{\linewidth}#1\end{minipage}\switchcolumn*}
\newcommand{\EndParallel}{\end{paracol}\end{spacing}}

% --- ZNAKI SPECJALNE ---
% 
%\newcommand †
% --------------------------

\begin{document}
\pagestyle{plain}% nagłówek i stopka pusta	
\setlength{\parindent}{0pt}
\gresetbracerendering{font}
\gresetinitiallines{0}\gabcsnippet{<c>7. </c>Ut(f) si(gh)ne(h) ti(h)mó(h)re,(h) de(h) ma(h)nu(h) in(h)i(h)mi(h)có(h)rum(h) nos(h)tró(h)<i><ul>rum</i></ul>(hg) <i><ul>li</i></ul>(ixgi)<i><ul>be</i></ul>(h hr[ocb:1{])<b>rá</b>(hg[ocb:0}])ti,(gh..) <c>*</c>(:) ser(h)vi(h)<i><ul>á</i></ul>(g)<i><ul>mus</i></ul>(f) <b>il</b>(gh gr)li.(g.) (::)}
\red{8.} In sanctitáte, et justíti\ul{\textit{a}} \ul{\textit{co}}\ul{\textit{ram}} \textbf{ip}so,~* ómnibus di\ul{\textit{é}}\ul{\textit{bus}} \textbf{nos}tris.\\
\red{9.} Et tu, puer, Prophéta Altís\ul{\textit{si}}\ul{\textit{mi}} \ul{\textit{vo}}\textbf{cá}\textbf{be}ris:~* præíbis enim ante fáciem Dómini, paráre \ul{\textit{vi}}\ul{\textit{as}} \textbf{e}jus:\\
\red{10.} Ad dandam sciéntiam salú\ul{\textit{tis}} \ul{\textit{ple}}\ul{\textit{bi}} \textbf{e}jus:~* in remissiónem peccató\ul{\textit{rum}} \ul{\textit{e}}\textbf{ó}rum:\\
\red{11.} Per víscera misericórdi\ul{\textit{æ}} \ul{\textit{De}}\ul{\textit{i}} \textbf{nos}tri:~* in quibus visitávit nos, óri\ul{\textit{ens}} \ul{\textit{ex}} \textbf{al}to:\\
\red{12.} Illumináre his, qui in ténebris, et in um\ul{\textit{bra}} \ul{\textit{mor}}\ul{\textit{tis}} \textbf{se}dent:~* ad dirigéndos pedes nostros in \ul{\textit{vi}}\ul{\textit{am}} \textbf{pa}cis.
\gresetinitiallines{1}\greannotation{Ad Bene-}\greannotation{dictus}\greannotation{1. g}\gabcsnippet{(c4) TRá(f)di(c')tor(d) au(dh)tem(ixh_i) () de(h')dit(g) e(e')is(f) si(g)gnum,(fe~) di(d.)cens :(d.) (;) Quem(d) o(fe)scu(d)lá(e')tus(f) fú(gh)e(g)ro,(f.) (,) i(f_g)pse(f) est,(c.) (;) te(e)né(g_[oh:h]e)te(f_e) e(d.)um.(d.) (::)}
\rubryka{Flexis genibus:}\relax
\gresetinitiallines{1}\greannotation{5.}\gabcsnippet{(c4) CHri(ffg)stus(f.) <c>*</c>(,) fa(fg)ctus(f) est(f') pro(f) no(ghF'ED.g_[oh:h]e/fg)bis(f.) (;) ob(f)é(hf/gh!jj/kjk)di(jij)ens(h_g) (,) us(h!jjh)que(f.) ad(ixfiH'Gh) mor(f_[oh:h]g_[oh:h]f_[oh:h])tem,(f.) (::)}
%% --- Pater noster ---
\hspace{5mm}Pater noster.~\textit{\red{secreto}}\\
\rubryka{Aliquantulum altius:}
\subsekcja{Psalm 50}
\red{1.} Miserére mei Deus, \red{*} secúndum magnam misericórdiam tuam.\hfill\break
\red{2.} Et secúndum multitúdinem miseratiónum tuárum, \red{*} dele iniquitátem meam.\hfill\break
\red{3.} Amplius lava me ab iniquitáte mea: \red{*} et a peccáto meo munda me.\hfill\break
\red{4.} Quóniam iniquitátem meam ego cognósco: \red{*} et peccátum meum contra me est semper.\hfill\break
\red{5.} Tibi soli peccávi, et malum coram te feci: \red{*} ut justificéris in sermónibus tuis, et vincas cum judicáris.\hfill\break
\red{6.} Ecce enim in iniquitátibus concéptus sum: \red{*} et in peccátis concépit me mater mea.\hfill\break
\red{7.} Ecce enim veritátem dilexísti: \red{*} incérta et occúlta sapiéntiæ tuæ manifestásti mihi.\hfill\break
\red{8.} Aspérges me hyssópo, et mundábor: \red{*} lavábis me, et super nivem dealbábor.\hfill\break
\red{9.} Audítui meo dabis gáudium et lætítiam: \red{*} et exsultábunt ossa humiliáta.\hfill\break
\red{10.} Avérte fáciem tuam a peccátis meis: \red{*} et omnes iniquitátes meas dele.\hfill\break
\red{11.} Cor mundum crea in me, Deus: \red{*} et spíritum rectum ínnova in viscéribus meis.\hfill\break
\red{12.} Ne projícias me a fácie tua: \red{*} et spíritum sanctum tuum ne áuferas a me.\hfill\break
\red{13.} Redde mihi lætítiam salutáris tui: \red{*} et spíritu principáli confírma me.\hfill\break
\red{14.} Docébo iníquos vias tuas: \red{*} et ímpii ad te converténtur.\hfill\break
\red{15.} Líbera me de sanguínibus, Deus, Deus salútis meæ: \red{*} et exsultábit lingua mea justítiam tuam.\hfill\break
\red{16.} Dómine, lábia mea apéries: \red{*} et os meum annuntiábit laudem tuam.\hfill\break
\red{17.} Quóniam si voluísses sacrifícium, dedíssem útique: \red{*} holocáustis non delectáberis.\hfill\break
\red{18.} Sacrifícium Deo spíritus contribulátus: \red{*} cor contrítum et humiliátum, Deus, non despícies.\hfill\break
\red{19.} Benígne fac, Dómine, in bona voluntáte tua Sion: \red{*} ut ædificéntur muri Jerúsalem.\hfill\break
\red{20.} Tunc acceptábis sacrifícium justítiæ, oblatiónes, et holocáusta: \red{*} tunc impónent super altáre tuum vítulos.
\subsekcja{Oratio}
\red{\hspace{5mm}\small\textit{Sine \textup{Oremus}, simili voce:}}
\parbox{\textwidth}{\inicjal{R}{éspice}, quǽsumus, Dómine, super hanc famíliam tuam, pro qua Dóminus noster Jesus Christus non dubitávit mánibus tradi nocéntium, et crucis subíre torméntum:}
\rubryka{Et sub silentio concluditur:}
\parbox{\textwidth}{\inicjal{Q}{ui} tecum vivit et regnat in unitáte Spíritus Sancti, Deus, per ómnia sǽcula sæculórum. \red{\Rbar.}~Amen.}
\end{document}